% ========== 第一章：绪论 ==========
\chapter{绪论}

\section{研究背景与动机}

近年来，以 Stable Diffusion~\cite{rombach2022latentdiffusion} 为代表的文生图（Text-to-Image）扩散模型在图像生成质量上取得了突破性进展。然而，当任务从"根据文本描述生成一张好看的图片"转变为"精确控制生成图像的相机视角与空间结构"时，现有模型暴露出明显局限——文本 prompt（如"俯拍""镜头向左移动"）本质上只能提供语义层面的模糊引导，模型难以精确理解相机的位置、角度、视场等几何信息，导致生成结果的视角不可控、结构不稳定\cite{zhang2023controlnet}。

与此同时，在电影摄影、3D 内容创作、虚拟现实等应用场景中，用户对多视角一致性和精确镜头控制的需求日益增长。例如，给定物体的正面参考图，用户期望生成该物体从侧面、顶部等目标视角观察时所见的外观。这类任务要求模型不仅理解图像内容，还需具备对三维空间变换的推理能力。NeRF~\cite{mildenhall2020nerf} 与 3D Gaussian Splatting~\cite{kerbl3dgaussians} 等方法可以从多视角图像集合中重建场景并合成新视角，但它们通常需要数十至上百张标定图像作为输入，难以从单张参考图直接泛化。

本项目的核心动机在于：\textbf{能否让扩散模型以单张参考图为输入，结合显式的相机旋转条件，直接生成目标视角下的、结构一致的多视角图像？} 为实现这一目标，我们将抽象的相机控制问题转化为扩散模型可理解的结构控制问题——通过深度图（Depth Map）与轮廓图（Canny Edge Map）作为中间条件表示，借助 ControlNet 的条件注入能力与 LoRA 的轻量化微调，构建一个面向 AI 图像生成镜头控制的多模态条件扩散原型系统。

\section{问题定义}

本项目的核心任务可形式化定义为：给定一张参考图像 $I_{\text{ref}}$ 与一个三维相机相对旋转向量 $\mathbf{r} = [\Delta\theta_{\text{azimuth}}, \Delta\theta_{\text{elevation}}, \Delta\theta_{\text{roll}}]^\top$（单位：度），生成目标视角下的结构图 $I_{\text{target}}$，使得 $I_{\text{target}}$ 在物体结构、轮廓与空间关系上与真实目标视角尽可能一致。系统采用 InstructPix2Pix~\cite{brooks2023instructpix2pix} 作为基础扩散模型，在其 UNet 的注意力层上挂载 LoRA 适配器，仅微调少量参数以适应这一新任务；推理阶段通过将大角度旋转分解为多个小步长、串联迭代生成的方式支持任意角度的视角变换。

\section{主要贡献}

本文的主要贡献包括：

\begin{enumerate}[label=(\arabic*)]
    \item \textbf{结构图视角变换方案}：提出以深度图/轮廓图作为中间表示，基于 InstructPix2Pix 原生 8 通道 UNet 进行 LoRA 微调，实现"参考结构图 + 旋转向量 → 目标结构图"的条件生成，避免了从零训练条件注入权重的开销。
    \item \textbf{RotationEncoder 旋转条件编码}：设计三维旋转向量的 sin/cos 编码方案与 MLP 映射模块，将其转换为可与 CLIP 文本嵌入拼接的伪文本 token，使扩散模型能够在 cross-attention 中感知相机位姿信息。
    \item \textbf{Blender 多视角数据集自动生成管线}：搭建了从场景构建、相机轨迹采样、批量渲染到深度图/轮廓图后处理的全自动化合成数据管线，可灵活扩展至各类 3D 资产。
    \item \textbf{多步迭代推理策略}：将大角度旋转分解为若干小角度步长的串联迭代，每步以上一步输出作为新的条件图，有效降低了单次推理的任务难度，使模型在训练阶段仅需学习小角度变换即可覆盖全范围旋转。
    \item \textbf{完整的训练—验证—部署流程}：提供了 LoRA 训练、收敛自动验证、Gradio 交互推理界面及 ComfyUI 工作流部署的全链路实现。
\end{enumerate}

\section{章节安排}

本文剩余部分组织如下：第二章综述扩散模型、可控生成、LoRA 微调、深度估计等相关工作；第三章详述整体方法设计，包括双分支架构、条件注入机制与多步推理策略；第四章介绍基于 Blender 的多视角数据集构建管线；第五章给出训练、收敛验证、推理应用及 ComfyUI 部署的实现细节；第六章分析实验结果与当前局限性；第七章总结全文并展望未来方向。
